\pageheading{EVIDENCE / FOUNDATIONS}{Established results and open questions}
The foundational families close selected structural questions and admit a reproducible visual treatment. Historical priority is a distinct question that the bounded literature review does not settle.

\begin{tabular}{@{}>{\raggedright\arraybackslash}P{0.46\linewidth}>{\raggedright\arraybackslash}P{0.47\linewidth}@{}}
\toprule
Established results & Remaining questions\\\midrule
Positive lower and finite upper dilation type & Existence of an optimal type limit and logarithmic modulation\\
Joint analytic parameter dependence and compact tap bounds & General minimal-bank identifiability\\
Exact arithmetic parameter recovery up to modulus radical & Stable recovery within bounded prime classes\\
All-prime signed-atom detectors and positivity domain & Sharp sieve-root asymptotics and multiplicities\\
Global orbit-zeta continuation and shifted determinants & Functional equations and derivative-weighted zero sums\\
Integer coefficient and exact polyphase/tap machinery & Historical priority of the candidate extensions\\
\bottomrule\end{tabular}

\claim{Structural distinctions among the families}
The asymmetric dilation family has intrinsic parameters, and its off-diagonal
members cannot be obtained by rescaling the diagonal family. Its nonnegative
real scale regimes are classified in Appendix~\ref{app:sequence}.
The arithmetic inverse and signed-resolution results address different
questions about parameter recovery and finite detection. The orbit
continuation applies established spectral methods to a specified seed,
with the branch conventions and hypotheses checked locally.

\claim{A common finite-resolution lesson}
A weak tap can change asymptotic order while preserving compact convergence.
A large excluded prime can change an exact arithmetic parameter while hiding
within a small bounded-window error. A negative atom can evade low-degree
positivity and become detectable at a calculable larger degree.
These are three different proved mechanisms. Their shared resolution theme
does not make them identical constructions.

\Needspace{13\baselineskip}
\pageheading{EVIDENCE / EXACT CHECKS}{Evidence with declared ranges}
The archived foundational review reconstructed five proof chapters and examined
their mathematical statements. The recorded release run executed ten scripts, with
25 pinned inputs and byte-identical isolated snapshots. All inputs remained
unchanged, every declared result was present, and the schema validator passed.

\[
\boxed{19{,}871\text{ exact assertion statements}\quad+
5{,}576\text{ numerical assertion statements passed}.}
\]
The exact total comprises 19,331 principal checks and 540 separately recorded
helper checks in the inverse reconstruction. Operational path and process
guards are excluded. Counts are assertion statements, not counts of theorems
or every scalar operation inside a grouped assertion.

In that recorded run, the API contributes 8,038 exact assertions and 18 catalog rows through index
12. Dilation type and continuity add 6,091 exact checks. Arithmetic inverse
examples include rank 10,000, negative depth 500, and large excluded primes.
The multiple-prime run tests 102 positive parameter/depth cases, 1,395
orthogonality conditions, 381 recurrence cases, and 40 conditional pencils.
Spectral continuation tests classical unshifted and shifted anchors.

\claim{Recovered failures are retained}
The preserved audit documents early environment and plotting issues, an unsuitable absolute determinant residual, and the previously known one-dilation case.
Two API validation defects were reproduced and fixed. A draft
sequence example had two miscopied integers and was corrected against the
exact catalog. Initial figure footer and label overlaps were repaired.
A direct publisher-source retrieval failed; an authenticated publisher
preview supplied the classical Laplace-uniqueness comparison instead.

Exact analytic tail envelopes are distinguished from floating-point rounding.
All decimal arithmetic and root approximations are ordinary high precision,
not interval certificates. No analytic Lean certificate is claimed.
The universal conclusions rely on the proofs, not on the number of checks.

\pageheading{EVIDENCE / REPRODUCIBILITY}{A reproducible research package}
The package preserves the research manuscripts, code, catalogs, review records
and bounded execution evidence. The typeset monograph uses 32 vector figures:
18 for the foundational families and 14 for coupled systems and finite
resolution. The corresponding generators and preview images remain in the
research archive. Input and output hashes identify the retained evidence.

The source files preserve the earlier claim ledger and open questions as provenance. Later proof chapters supersede the earlier statements where a new argument supplies a closed result. The distinction is recorded in the verification documents rather than folded into the mathematical theorem count.

\claim{Reproduce the work}
Build the book with \texttt{python monograph/build.py} from the package root.
The script requires XeLaTeX and resolves contents, citations and cross-references
through successive passes. Research commands are run from the
\texttt{research/} directory, whose \texttt{requirements.txt} records the
scientific dependencies. The package README supplies complete commands for
the isolated foundational and coupled verification suites. Figure generators
rebuild the atlas from formulas and recorded numerical examples.

The standard-library API is portable. Numerical checks additionally require
mpmath. No network access, external submission, or dependency installation is
part of the verification run. PDF byte identity can change with timestamps
and font environments; mathematical outputs retain their own provenance.

The remaining questions admit concrete starting points: type-limit analysis
for dilation, precision-controlled arithmetic inversion within bounded
parameter classes, sharper sieve-root asymptotics, and spectral weighted sums.
They remain research questions rather than hidden dependencies of the
proved statements in this volume.
